\section{Basis Encoding of Functions}

Let us now restrict to relations, which have an expression by functions.
We in this section then show, how contractions of their encodings can be exploited in function evaluation.

\subsection{Definition}

%We now generalize the representation scheme towards maps between arbitrary unstructured sets.

\begin{definition}[Basis encoding of maps]
    \label{def:functionRelationEncoding}
    Any map
    \begin{align*}
        \exfunction\defcols \inset \rightarrow \outset
    \end{align*}
    can be represented by a relation
    \begin{align*}
        \exrelationof{\exfunction} \coloneqq \left\{ \big(x,\exfunction(x)\big) \wcols x \in\inset \right\} \subset \inset \times \outset \, .
    \end{align*}
    Given a enumeration of the sets by $\indvariableof{\insymbol}$ and $\indvariableof{\outsymbol}$ we define the basis encoding of $\exfunction$ as the tensor
    \begin{align*}
        \bencodingofat{\exfunction}{\indvariableof{\outsymbol},\indvariableof{\insymbol}}
        = \onehotmapofat{\exrelationof{\exfunction}}{\indvariableof{\insymbol},\indvariableof{\outsymbol}}  \, .
    \end{align*}
\end{definition}

\begin{remark}[Reduction to images]
    % Image enumeration
    When $\exfunction$ maps into a set of infinite cardinality, we restrict $\outset$ to the image of $\exfunction$ and enumerate the image by a variable $\indvariableof{\exfunction}$.
    This scheme is applied, when $\exfunction$ is itself a tensor, i.e. $\outset=\rr$.
    While the variable $\indvariableof{\exfunction}$ can in general be of the same cardinality as the domain set $\inset$, it will be valued in $[2]$ when considering boolean tensors.
\end{remark}

% Characterization of the directed and boolean tensors
We notice that any basis representation of a function is also a directed tensor with incoming variables to the domain and outgoing variables to the image.
It furthermore holds, that the set of directed and boolean tensors is characterized by the basis encoding of functions.
This is shown in the next theorem, by the claim that any boolean tensor which is directed is the basis representation of a function.

\begin{theorem}
    \label{the:bencodingDirected}
    Let $\inset,\outset$ be sets and $\exrelation\subset\inset\times\outset$ a relation.
    If and only if there exists a map $\exfunction:\inset\rightarrow\outset$ such that $\exrelation=\exrelationof{\exfunction}$, the basis encoding $\bencodingof{\exfunction}$ is a directed tensor with $\indvariableof{\insymbol}$ incoming and $\indvariableof{\outsymbol}$ outgoing.
\end{theorem}
\begin{proof}
    \proofrightsymbol{}:
    When $\exfunction$ is a function, we have for any $\indindexofin{\insymbol}$
    \begin{align*}
        \sum_{\indindexofin{\outsymbol}} \bencodingofat{\exfunction}{\indexedindvariableof{\outsymbol},\indexedindvariableof{\insymbol}}
        =  \bencodingofat{\exfunction}{\indvariableof{\outsymbol}=\invindexinterpretationofat{\outsymbol}{\exfunctionat{\indexinterpretationofat{\insymbol}{\indindexof{\insymbol}}}},\indexedindvariableof{\insymbol}}
        = 1 \, .
    \end{align*}
    Thus, $\bencodingofat{\exfunction}{\indvariableof{\outsymbol},\indvariableof{\insymbol}}$ is a directed tensor with variables $\indvariableof{\insymbol}$ incoming and $\indvariableof{\outsymbol}$ outgoing.

    \proofleftsymbol{}:
    Conversely let there be a relation $\exrelation$, such that $\bencodingof{\exrelation}$ is directed.
    To this end, we observe that for any $\indindexofin{\insymbol}$ the tensor
    \begin{align*}
        \onehotmapofat{\exrelation}{\indexedindvariableof{\insymbol},\indvariableof{\outsymbol}}
    \end{align*}
    is a boolean tensor with coordinate sum one and therefore a basis vector.
    It follows that the function $\exfunction : \inset \rightarrow \outset $ defined for $x\in\inset$ as
    \begin{align*}
        \exfunctionat{x}
        = \indexinterpretationofat{\outsymbol}{\invonehotmapof{\onehotmapofat{\exrelation}{\indvariableof{\insymbol}=\invindexinterpretationofat{\insymbol}{x},\indvariableof{\outsymbol}}}}
    \end{align*}
    is well-defined.
    We then have by construction
    \begin{align*}
        \bencodingofat{\exfunction}{\indvariableof{\outsymbol},\indvariableof{\insymbol}}
        & = \sum_{\indindexofin{\insymbol}}
        \onehotmapofat{\exfunction(\indindexof{\insymbol})}{\indvariableof{\outsymbol}} \otimes
        \onehotmapofat{\indindexof{\insymbol}}{\indvariableof{\insymbol}} \\
        & =  \sum_{\indindexofin{\insymbol}} \onehotmapofat{\exrelation}{\indexedindvariableof{\insymbol},\indvariableof{\outsymbol}} \otimes
        \onehotmapofat{\indindexof{\insymbol}}{\indvariableof{\insymbol}} \\
        & = \onehotmapofat{\exrelation}{\indvariableof{\outsymbol},\indvariableof{\insymbol}}
    \end{align*}
    and therefore by \defref{def:functionRelationEncoding} $\exrelation=\exrelationof{\exfunction}$.
\end{proof}

% Grid sets
We are specially interested in sets of states of a factored system, which amounts to the case in \defref{def:functionRepresentation}.
Those state sets have a decomposition into a cartesian product of $\atomorder$ sets
\[ \arbset = \facstates \, . \]
The most obvious enumeration of the set $\arbset$ is therefore by the collection of state variables $\{\catvariableof{\atomenumerator}\wcols \atomenumeratorin \}$.
Functions between states of factored systems with $\atomorder_{\insymbol}$ and $\atomorder_{\outsymbol}$ state variables can be represented by $\atomorder_{\insymbol}+\atomorder_{\outsymbol}$-ary relations and \defref{def:functionRelationEncoding} has an obvious generalization to this case with multiple enumeration variables.

%% NOT NEEDED -> Done in propositional logics
%% Conditional
%Since the basis encoding of any map between factored systems is directed, it can be interpreted by a conditional probability tensor, as we state next.
%
%%% Maps
%\begin{corollary}%\label{the:condProbFunctionRepresentation}
%	The basis encoding $\bencodingof{\exfunction}$ (see \defref{def:functionRepresentation}) of a function $\exfunction$ between factored systems is a conditional probability tensor, where the legs to the image system are the conditions and the legs to the target system the distribution legs.
%\end{corollary}
%
%%% Deterministic by construction
%These are deterministic conditional probability tensors, in the sense that any slice with respect to the input variables is a basis tensor.
%Through contractions with distribution tensors (e.g. distributions in domain systems) they get stochastic.
%This is for example the case in the empirical distribution, which can be understood as the forwarding of the uniform distribution on the sample enumeration.



\subsection{Function Evaluation}

We now justify the nomenclature of basis calculus, by showing that contraction with basis elements produce the one-hot encoded function evaluation.

\begin{theorem}[Function evaluation in Basis Calculus]
    \label{the:basisCalculus}
    Retrieving the value of the function $\exfunction$ at a specific state is then the contraction of the tensor representation with the one-hot encoded state.
    For any $\arbelement\in\inset$ we have
    \begin{align*}
        \onehotmapofat{\invindexinterpretationofat{\outsymbol}{\exfunctionat{\arbelement}}}{\indvariableof{\outsymbol}}
        = \contractionof{
            \bencodingofat{\exformula}{\indvariableof{\outsymbol},\indvariableof{\insymbol}},
            \onehotmapofat{\invindexinterpretationofat{\insymbol}{\arbelement}}{\indvariableof{\insymbol}}
        }{\indvariableof{\outsymbol}} \, .
    \end{align*}
    Thus, we can retrieve the function evaluation by the inverse one-hot mapping as
    \begin{align*}
        \exfunctionat{\arbelement} = \invonehotmapof{\contractionof{
            \bencodingofat{\exformula}{\indvariableof{\outsymbol},\indvariableof{\insymbol}},
            \onehotmapofat{\invindexinterpretationofat{\insymbol}{\arbelement}}{\indvariableof{\insymbol}}
        }{\indvariableof{\outsymbol}}} \, .
    \end{align*}
\end{theorem}
\begin{proof}
    From the representation
    \begin{align*}
        \bencodingofat{\exfunction}{\indvariableof{\outsymbol},\indvariableof{\insymbol}}
        & =  \sum_{\indindexofin{\insymbol}}
        \onehotmapofat{(\invindexinterpretationof{\outsymbol} \circ \exfunction \circ \indexinterpretationof{\insymbol}) \indindexof{\insymbol}
        }{\indvariableof{\insymbol}}
        \otimes
        \onehotmapofat{\indindexof{\insymbol}}{\indvariableof{\insymbol}}
    \end{align*}
    and the orthonormality of the one-hot encodings of the input enumeration we get
    \begin{align*}
        \contractionof{
            \bencodingofat{\exformula}{\indvariableof{\outsymbol},\indvariableof{\insymbol}},
            \onehotmapofat{\invindexinterpretationofat{\insymbol}{\arbelement}}{\indvariableof{\insymbol}}
        }{\indvariableof{\outsymbol}}
        = \onehotmapofat{\invindexinterpretationofat{\outsymbol}{\exfunctionat{\arbelement}}}{\indvariableof{\outsymbol}} \, . & \qedhere
    \end{align*}
\end{proof}

% Comparsion with coordinate calculus
In comparison with the Coordinate Calculus scheme (see \theref{the:coordinateCalculus}), the Basis Calculus produces basis vectors of a functions evaluation instead of scalars.
While this seems to produce unnecessary redundancy in representing a function, we will see in the following section, that this scheme is efficient in representing compositions of functions.
